The forecaster is not a standalone artefact. It sits inside a control loop in which its output drives battery dispatch and curtailment on a photovoltaic micro-grid, so a forecasting failure can become a control failure, and the property tested in this report has to be read as a security property and not only as an accuracy one. The NIST AI Risk Management Framework\autocite{tabassiArtificialIntelligenceRisk2023} is the reference adopted for that reading, because it separates the trustworthiness characteristics evidenced here, validity and reliability, security and resilience, from the governance functions that turn evidence into a control, and because it asks that known limitations be mapped, measured and managed rather than merely disclosed.

\subsection{The limitation is a weakness before it is a vulnerability}

Characterising the limitation correctly determines everything that follows. Phase-locking is deterministic, attacker-independent and present at deployment time. It is not introduced by an adversary and it is not a defect in the training data; it is a consequence of the tokenisation geometry, and it is closed-form. The candidate set follows from $(f_s,P,S)$ alone, and those three numbers are configuration, not weights. Structural aliasing is the phenomenon that may then appear in the learned representation, and phase-locking is necessary for it but not sufficient, so the candidate set bounds where the model can be blind while \nameref{sec:seven} is what establishes where it is.

An adversary therefore needs no model access, no gradients, no query budget and no knowledge of the training corpus to enumerate that set, and only observation of the deployed forecaster to find which of its members carry a deficit. In the taxonomy of adversarial machine learning\autocite{vassilevAdversarialMachineLearning2024} this places the weakness with the availability violations rather than with integrity or privacy, and it places the attack outside the white-box assumption, which is what a conventional threat model tends to miss because it reasons about perturbations of the input tensor rather than of the physical process.

A weakness becomes a vulnerability only when someone can reach it. For the leading threat model, adversarial masking, the five elements are the following. \textbf{Adversary capabilities}: no model access of any kind and no knowledge of the training corpus, only the ability to drive a periodic load on the plant, which an inverter switching schedule or a switchable load already provides. \textbf{Attack surface}: the telemetry stream upstream of tokenisation, reached by acting on the physical process rather than on the data path. \textbf{Manipulation of telemetry}: a small periodic component at a candidate frequency, which is legitimate plant behaviour and survives any plausibility check on the raw signal. \textbf{Operational consequence}: the component is present in the plant and under-represented in the model, so dispatch and curtailment are decided on an incomplete picture, and a detector reading the forecast residual is looking through the same tokeniser that lost it. \textbf{Residual risk}: overlap thins the stride branch alone, so the patch branch remains at every deployed configuration.

\begin{table*}[!t]
    \centering
    \small
    \setlength{\tabcolsep}{2.25pt}
    \begin{tabular*}{\textwidth}{@{\extracolsep{\fill}}>{\raggedright\arraybackslash}m{2.50cm}>{\raggedright\arraybackslash}m{3.00cm}>{\raggedright\arraybackslash}m{1.70cm}>{\raggedright\arraybackslash}m{3.85cm}>{\raggedright\arraybackslash}m{3.25cm}>{\centering\arraybackslash}m{1.45cm}@{}}
        \toprule
        \textbf{Risk} & \textbf{Description} & \textbf{AI RMF} & \textbf{Evidence in this report} & \textbf{Control available} & \textbf{Residual} \\
        \midrule
        Adversarial masking & A periodic component injected on a candidate frequency is under-represented and escapes forecast-residual detection & \shortstack[l]{MAP\\MEASURE} & Precondition is $e^{\theta_{\mathrm{lock}}}>1$ (Model~B) with the site locations of $\kappa_S$ (Model~D2); the pathway itself needs the paired recovery ratio of Model~A. Not yet evaluated & Spectral monitor on the raw stream, upstream of tokenisation, at the enumerated candidate set & high \\
        \cmidrule(lr){1-6}
        Systematic blind spot with no adversary & A legitimate plant harmonic sits on a candidate frequency and is persistently under-represented & \shortstack[l]{MAP\\MEASURE\\MANAGE} & Candidate set is closed-form in $(f_s,P,S)$ and needs no fit; the cost at those sites is decided by $e^{\theta_{\mathrm{lock}}}$. Not yet evaluated & Choose $(P,S)$ against the known spectral content of the plant; \valueof{PartiallyObservable} degradation & medium \\
        \cmidrule(lr){1-6}
        Configuration drift & A change of sampling interval, patch size or checkpoint silently relocates the blind set & \shortstack[l]{GOVERN\\MAP} & $H3a$: relocation is predicted in closed form by $f_s/S$, and whether the sites follow it is decided by $\kappa_S$. Not yet evaluated & Recompute the candidate set from the knowledge base at every configuration change & low \\
        \cmidrule(lr){1-6}
        False assurance & Overlap is increased and treated as a fix & \shortstack[l]{GOVERN\\MANAGE} & $M1$ is decided by $\delta_O$ under a two-leg rule and neither leg is evaluated. Independently of any fit, the patch branch is invariant to $S$ by construction & Do not credit overlap as a control; record the unestablished status in the model card & high \\
        \cmidrule(lr){1-6}
        Defence placed inside the blind spot & Phase randomisation, or anomaly detection on forecast residuals & MANAGE & $H2$ is decided by $\sigma_\phi$ (Model~C): a deficit independent of phase makes randomising it useless. Not yet evaluated & Move detection upstream of tokenisation; treat residual-based detection as insufficient alone & medium \\
        \bottomrule
    \end{tabular*}
    \caption{AI RMF names the function of the framework\autocite{tabassiArtificialIntelligenceRisk2023} under which each item is handled. Residual is the exposure that survives the control in the adjacent column, judged qualitatively: high where no control available to an operator closes the exposure, medium where a control exists but is partial or requires design-time choices, low where the exposure is closed by a procedure that follows from the arithmetic. Each entry names the estimand specified in \autoref{sec:six} that decides it; no posterior is asserted here, since the fits are not yet available. The two entries carried at high residual are those for which no control available to an operator closes the exposure: Adversarial masking, because a detector placed after tokenisation inherits the blind spot it is meant to find, and False assurance, because the pre-registered mitigation cannot reach the patch branch by the arithmetic of \autoref{sec:eight} whatever the fits return.}
    \label{tab:riskRegister}
\end{table*}


\subsection{What the evidence supports and what it does not}

The results of \nameref{sec:seven} sharpen this picture in one direction and weaken it in another, and reporting both is the point of having measured rather than asserted.

Three findings strengthen it. The representational cost at a candidate frequency is confirmed, at a $27\%$ codelength expansion with a credible interval well clear of the threshold, so the precondition of every scenario below is established rather than assumed. The sites track their own arithmetic exactly, $\kappa_S=1.000\,[0.992,1.008]$, so an adversary or an auditor can compute the target set in closed form from configuration alone and be right about where it sits. And the deficit does not depend on the phase of the signal, which removes any need for the adversary to control alignment and simultaneously removes phase randomisation from the operator's list of defences.

One finding weakens it, and it is the most consequential. The behavioural half of $H1$ came out inconclusive and, worse, prior-dependent: the forecast is not demonstrably attenuated at a candidate frequency. The strongest form of the threat model, in which the injected component is simply absent from the forecast, is therefore \textbf{not} supported by this evidence. What is supported is the weaker and still security-relevant statement that the internal representation carries a measurable, predictable and configuration-derived cost at an enumerable set of frequencies, while the decoder appears able to compensate for it under the conditions tested. A representation that is degraded but compensated is a brittleness rather than a demonstrated exploit: it is a margin that a distribution shift, a longer horizon or a lower-amplitude component could consume, and none of those were tested.

$T1$ in \autoref{tab:riskRegister} is therefore carried as a plausible but unproven pathway with a confirmed precondition, not as a demonstrated attack. What would settle it is specific: the same paired contrast at the amplitude, horizon and sampling interval of the deployed telemetry rather than of the probe, on a model that meets the convergence thresholds.

\subsection{Threat models beyond the adversarial one}

The framework asks for systematic vulnerability, not only for attacks, and the more likely exposure has no adversary in it. On one-minute telemetry, the arithmetic of \nameref{sec:two} puts the patch branch of a $P{=}16$ tokeniser at periods of $16/k$ minutes, that is at $16$, $8$, $5.3$, $4$, $3.2$, $2.7$ and $2.3$ minutes. Those are the timescales of inverter duty cycling, of maximum-power-point dither and of cloud-shadow transients: ordinary plant behaviour, not an attack, sitting on frequencies the tokeniser is least able to represent. A monitoring system built on such a forecaster would be systematically less sensitive there, permanently and for reasons no operational log would reveal. That is $T2$, and it is why the candidate set belongs in the deployment documentation of such a system.

The third exposure is governance rather than signal processing. Because the blind set is a function of configuration, a change that no one would route through a security review, a resampling of telemetry from one minute to thirty seconds, a swap of checkpoint, a patch size chosen for latency, relocates it. $H3a$ is what makes this tractable: the relocation is not merely predictable in principle but confirmed to follow $f_s/S$ one for one, so recomputing the set is arithmetic rather than re-measurement.

\subsection{Controls, residual exposure and conclusion}

Two controls in this design are real. Detection has to be placed before tokenisation, on the raw telemetry, because anything placed after it inherits the blind spot it is meant to guard; a spectral monitor on the raw stream, evaluated on the enumerated candidate set, is inexpensive precisely because that set is closed-form. And the semantic layer of \autoref{sec:appendixA} degrades actuation rather than failing silently: an assessment that is merely \valueof{PartiallyObservable} reduces the battery cap and lowers decision confidence, so uncertainty about the representation reaches the actuator instead of being resolved by assumption. That is the resilience the framework asks for, and it is the part of the design that does not depend on the disputed behavioural claim.

One control cannot be credited, and the report is more useful for saying so. The pre-registered mitigation of \nameref{sec:eight}, increased patch overlap, is not established: the posterior for $\delta_O$ points in the predicted direction but the fit it is drawn from fails its convergence checks and the second leg of the rule is not met, so the claim is inconclusive rather than confirmed. The arithmetic bounds it in any case, since overlap moves the stride branch and leaves the patch branch untouched, so even a mitigation established on this design would be partial by construction. An operator who raises the overlap ratio and records the risk as mitigated has bought a false assurance, which is $T4$ and the reason it is carried at high residual exposure.

The honest conclusion is therefore narrower than the one this section could have asserted and better supported than one it could have assumed. Structural aliasing in a patch-based forecaster is a configuration-determined, attacker-independent weakness whose location is now measured rather than conjectured; its representational cost is established, its propagation to the forecast is not, and the mitigation carried through the design is neither established nor capable of closing it. The residual exposure is significant, it is concentrated in the patch branch, and it cannot be managed inside the model. It has to be managed around it: by choosing the geometry against the known spectral content of the plant, by monitoring the raw signal upstream of the tokeniser, and by treating the candidate set as a deployment parameter that is recomputed whenever the configuration it derives from changes.
